% The 24-Meet of Lo Shu and Durer --- paper source
%
% Public-facing claims are controlled by docs/CLAIM_CROSSWALK.md.
% Analytic statements are proved in the text.  Finite enumerations and exact
% replay statements are labelled as Certified Propositions and point to
% repository replay artifacts.

\documentclass[11pt,a4paper]{article}

\usepackage[T1]{fontenc}
\usepackage[utf8]{inputenc}
\usepackage{lmodern}
\usepackage{amsmath,amssymb,amsthm,mathtools}
\usepackage{booktabs}
\usepackage{array}
\usepackage{enumitem}
\usepackage{microtype}
\usepackage[margin=1in]{geometry}
\usepackage{xcolor}
\usepackage{xurl}
\usepackage{xspace}
\usepackage{hyperref}

\hypersetup{
  colorlinks=true,
  linkcolor=blue!55!black,
  citecolor=blue!55!black,
  urlcolor=blue!55!black,
  pdftitle={The 24-Meet of Lo Shu and Durer: Bounded Magic Spectra, Sagrada Rays, and F2-to-the-4 Transport},
  pdfauthor={Oleksiy Babanskyy},
  pdfsubject={Bounded magic-square endpoint certificates},
  pdfkeywords={magic squares, Lo Shu, Durer square, Sagrada Familia, F2-to-the-4, affine geometry, finite certificates}
}
\emergencystretch=2em

% ---- Theorem environments ----
\newtheorem{theorem}{Theorem}[section]
\newtheorem{proposition}[theorem]{Proposition}
\newtheorem{certprop}[theorem]{Certified Proposition}
\newtheorem{lemma}[theorem]{Lemma}
\newtheorem{corollary}[theorem]{Corollary}
\newtheorem{conjecture}[theorem]{Conjecture}
\theoremstyle{definition}
\newtheorem{definition}[theorem]{Definition}
\theoremstyle{remark}
\newtheorem{remark}[theorem]{Remark}

% ---- Notation ----
\newcommand{\B}{\mathcal B}
\newcommand{\Spec}{\operatorname{Spec}}
\newcommand{\Ray}{\operatorname{Ray}}
\newcommand{\Meet}{\operatorname{Meet}}
\newcommand{\conv}{\operatorname{conv}}
\newcommand{\rank}{\operatorname{rank}}
\newcommand{\Aut}{\operatorname{Aut}}
\newcommand{\SNF}{\operatorname{SNF}}
\newcommand{\F}{\mathbb F}
\newcommand{\Z}{\mathbb Z}
\newcommand{\Q}{\mathbb Q}
\newcommand{\AGL}{\operatorname{AGL}}
\newcommand{\GL}{\operatorname{GL}}
\newcommand{\APD}{\operatorname{APD}}
\newcommand{\Durer}{D\"urer\xspace}

\title{The 24-Meet of Lo Shu and \Durer:\\
Bounded Magic Spectra, Sagrada Rays, and \(F_2^4\) Transport}
\author{Oleksiy Babanskyy}
\date{}

\begin{document}
\maketitle

\begin{abstract}
We isolate a bounded endpoint shared by two different magic-square
mechanisms.  On the \(3\times3\) side, repeated-entry magic squares with
entries in \([1,9]\) have an upper non-degenerate Lo Shu fiber at sum \(24\)
and a degenerate all-9 fiber at sum \(27\).  On the \(4\times4\) side, a
one-incidence Sagrada mask on the \Durer-complement square defines a bounded
ray \(D(t)=D-tM\) whose terminal point occurs at \(t=10\), again at sum
\(24\).  The weak meet is \(\{24,27\}\); after imposing non-degeneration on
the Lo Shu side and terminality on the \Durer/Sagrada side, the unique strong
meet is \(\{24\}\).

The endpoint also carries finite structure on the \Durer side.  The target
permutation diagonals drop from a dihedral subgroup \(D_4\le S_4\) to the
Klein four subgroup \(V_4\), and the \Durer value-minus-one cell model gives
an \(F_2^4\) affine transport layer: \(52\) affine source quaternes reduce to
\(36\) affine terminal quaternes.  Within this layer, a scoped parity-plane
mechanism explains the endpoint-24 bit-balance as a zero-residual affine
phenomenon; its final residual obstruction is a six-value finite check, not a
table-free symbolic theorem.  A broader order-four audit shows that
terminal \(24\) is not unique: it occurs in \(236\) square-mask pairs.  Inside
that landscape, the most stable finite class detected by the atlas is a
\(144\)-pair exact \(V_4\) subclass characterized by global
\(F_2^4\)-affineness; selected-affine near misses are separated inside the
main \(176\)-pair class by a quotient-shadow obstruction.  Thus \(24\) is not
a universal invariant of magic squares.  It is a structured endpoint in a
specific bounded meet, supported by exact finite certificates.
\end{abstract}

\noindent\textbf{Keywords.} Magic squares, Lo Shu, \Durer square, Sagrada
Familia square, bounded lattice points, permutation diagonals, Klein four
group, affine geometry over \(F_2\), finite certificates.

\smallskip
\noindent\textbf{MSC 2020.} 05B15 (primary), 05A15, 05E18, 11P21, 52B20,
20B25.

\section{Introduction}

The classical Lo Shu square and \Durer's order-four magic square belong to
different corners of the magic-square literature.  \Durer's square has a rich
record of symmetries, quaterne patterns, and Sagrada Familia variants
\cite{mathworldDurer,maritzSagrada,trenklerMelancholic,sudberyTesseract}.
Order-four magic squares also have a classical finite classification: up to
the standard equivalences there are \(880\) normal representatives
\cite{ollerenshawBondi}.  The present paper adds no universal numerological
slogan to this literature.  It puts one observation into a bounded,
certificate-controlled category.

The observation is that the value \(24\) occurs as the endpoint of two
unrelated-looking bounded mechanisms:
\begin{itemize}[leftmargin=*]
\item a non-degenerate upper fiber of repeated-entry \(3\times3\) magic
  squares with values in \([1,9]\);
\item the terminal point of a one-incidence Sagrada perturbation ray inside a
  \(4\times4\) \Durer-complement square with values in \([1,16]\).
\end{itemize}
The first mechanism is a lattice diamond.  The second is a bounded ray.  Their
weak numerical intersection contains \(24\) and \(27\), but \(27\) fails both
endpoint tests: it is the degenerate all-9 Lo Shu fiber and an interior point
of the \Durer/Sagrada ray.  The strong meet is therefore \(24\).

The manuscript has two layers of results.
\begin{enumerate}[label=\textbf{(\alph*)},leftmargin=*]
\item The central bounded-meet theorem is a short mathematical consequence of
  the two certified spectra.
\item The \Durer-side structure is an exact finite replay: quaterne transport,
  the \(D_4\to V_4\) permutation-diagonal drop, the \(F_2^4\) affine transport
  layer, the bounded-polytope guardrail, and the order-four endpoint atlas.
\end{enumerate}
We use an explicit claim-level convention: replay-backed statements are
labelled as \emph{Certified Propositions}.  These are not external theorems;
they are finite exact replays or finite enumerations whose scope is explicitly
bounded by repository artifacts, tests, and a claim crosswalk.  The controlling
ledger for this manuscript is
\[
\texttt{docs/CLAIM\_CROSSWALK.md},
\]
and the main certificate pack is
\[
\texttt{results/magic24\_certificate\_pack.json}.
\]

\paragraph{What is not claimed.}
The guardrails are part of the result.  The paper does not claim that \(24\)
is a universal invariant of magic squares, that the terminal square \(D(10)\)
is more perfect than \Durer's source square, or that \(D(10)\) is a vertex of
the bounded order-four magic polytope.  It also does not identify the
terminal \(V_4\) diagonal subgroup with a type-A poset cone, and it does not
claim that the \(F_2^4\) layer isolates only the historical \Durer/Sagrada cell
in the full order-four atlas.  These stronger statements are listed as
blocked claims in the reproducibility package.

\paragraph{Organization.}
Sections~\ref{sec:bounded}--\ref{sec:meet} prove the strong bounded meet.
Sections~\ref{sec:quaternes}--\ref{sec:f2} give the \Durer-side enrichment:
quaterne transport, permutation diagonals, and the \(F_2^4\) tesseract layer.
Section~\ref{sec:parity-plane} records the scoped parity-plane mechanism
behind the endpoint-24 bit-balance.  Section~\ref{sec:polytope} records the
bounded-polytope guardrail, and
Section~\ref{sec:order4} summarizes the order-four extension.  The appendices
collect the endpoint atlas, exact canonical \(V_4\) subclass, secondary
audits, and open problems.

\section{Bounded Magic Squares and One-Incidence Rays}
\label{sec:bounded}

\begin{definition}[Bounded repeated-entry magic squares]
For positive integers \(n,m,S\), let \(\B(n,m,S)\) be the set of integer
\(n\times n\) arrays \(A=(a_{ij})\) with entries in \([1,m]\) such that all
rows, all columns, and the two main diagonals sum to \(S\).  Repeated entries
are allowed.
\end{definition}

The repeated-entry convention is essential.  It lets the Lo Shu side form a
short bounded spectrum and lets the \Durer side move along a perturbation ray.
Thus the phrase ``Lo Shu side'' means the affine \(3\times3\) magic-square
family normalized around the Lo Shu center, not normal \(3\times3\) squares
with distinct symbols \(1,\ldots,9\).  It is separate from the distinct-entry
inside-out enumerations of magic and magilatin labellings
\cite{beckZaslavskyInsideOut,beckZaslavskyMagicLabellings,beckVanHerick}.
Throughout this paper, ``magic'' means row sums, column sums, and the two main
diagonal sums; panmagic or toroidal diagonal constraints are not part of the
ambient category unless explicitly stated.

\begin{definition}[One-incidence mask]
Let \(Q\) be an \(n\times n\) magic square.  A zero-one mask \(M\) is
one-incidence if every row, every column, and both main diagonals of \(M\)
have sum \(1\).  Then
\[
Q(t)=Q-tM
\]
is again magic, with magic sum decreased by \(t\), for every integer \(t\).
The bounded terminal time is the largest \(t\ge 0\) for which \(Q(t)\) remains
inside the prescribed entry interval.
\end{definition}

For \(4\times4\) squares, one-incidence masks are indexed by the
permutations \(\sigma\in S_4\) that use one cell in each row and column and
also one cell in each main diagonal.

\section{The Lo Shu Lattice Diamond}
\label{sec:loshu}

Every \(3\times3\) magic square can be parametrized as
\[
\begin{pmatrix}
g+a & g-a-b & g+b\\
g-a+b & g & g+a-b\\
g-b & g+a+b & g-a
\end{pmatrix},
\qquad S=3g.
\]
Imposing entries in \([1,9]\) gives a lattice diamond in the \((a,b)\)-plane:
\[
|a+b|\le m,\qquad |a-b|\le m,\qquad
m=\min(g-1,9-g).
\]
The diamond has
\[
2m^2+2m+1
\]
integer points.

\begin{certprop}[Bounded Lo Shu spectrum]
\label{prop:loshu-spectrum}
For \(S=3,6,\ldots,27\), the sizes of \(\B(3,9,S)\) are
\[
1,\;5,\;13,\;25,\;41,\;25,\;13,\;5,\;1.
\]
In particular,
\[
\Spec^+_{3,9}(15)=\{18,21,24,27\}.
\]
At \(S=24\) the fiber has \(5\) points, while at \(S=27\) the fiber is the
unique all-9 square.
\end{certprop}

\begin{proof}[Certification]
This is Claim IDs M24-C01--M24-C03 and M24-C57 in
\path{docs/CLAIM_CROSSWALK.md}.  The replay is in the certificate pack and
the Lo Shu lattice-polygon note:
\begin{itemize}[leftmargin=*]
\item \path{results/magic24_certificate_pack.json}
\item \path{docs/LO_SHU_LATTICE_POLYGON.md}
\end{itemize}
The formula above gives the closed-form count.
\end{proof}

\begin{table}[htbp]
\centering
\begin{tabular}{rrrrrrrrrr}
\toprule
\(S\) & 3 & 6 & 9 & 12 & 15 & 18 & 21 & 24 & 27\\
\midrule
\(|\B(3,9,S)|\) & 1 & 5 & 13 & 25 & 41 & 25 & 13 & 5 & 1\\
\bottomrule
\end{tabular}
\caption{Bounded repeated-entry \(3\times3\) magic-square counts with entries
in \([1,9]\).}
\label{tab:loshu-counts}
\end{table}

At \(S=24\), \(g=8\), \(m=1\), and the five parameter points are
\[
(0,0),\quad (1,0),\quad (0,1),\quad (-1,0),\quad (0,-1).
\]
At \(S=27\), \(g=9\), \(m=0\), and the diamond collapses to a single point.
This collapse is the non-degeneration distinction that will remove \(27\)
from the strong meet.

\section{The \Durer/Sagrada Ray}
\label{sec:durer-ray}

We use the \Durer-complement square
\[
D=
\begin{pmatrix}
1&14&15&4\\
12&7&6&9\\
8&11&10&5\\
13&2&3&16
\end{pmatrix},
\qquad S_D=34.
\]
This is the complement orientation of the classical \Durer square, obtained
by \(x\mapsto17-x\), chosen because it matches the Sagrada mask convention
used here.
Let \(M\) be the one-incidence mask indexed by \(2013\), selecting the four
cells
\[
(0,2),\quad (1,0),\quad (2,1),\quad (3,3).
\]
These cells contain \(15,12,11,16\).

\begin{figure}[htbp]
\centering
\[
\begin{pmatrix}
1&14&\boxed{15}&4\\
\boxed{12}&7&6&9\\
8&\boxed{11}&10&5\\
13&2&3&\boxed{16}
\end{pmatrix}
\]
\caption{The Sagrada mask \(2013\) on the \Durer-complement square.}
\label{fig:sagrada-mask}
\end{figure}

\begin{certprop}[Sagrada terminality]
\label{prop:sagrada-terminal}
The ray \(D(t)=D-tM\) remains inside \([1,16]\) exactly for
\[
0\le t\le 10.
\]
Its terminal magic sum is \(34-10=24\).  Among the eight one-incidence \Durer
masks, \(2013\) is the unique mask with terminal sum \(24\).  Equivalently,
it is the unique maximin mask: it uniquely maximizes the smallest selected
entry, hence uniquely maximizes the allowed downward time \(t_{\max}\).
\end{certprop}

\begin{proof}[Certification]
For selected values \(v_1,\ldots,v_4\), the bound \(D-tM\in[1,16]\) gives
\[
t_{\max}=\min_i(v_i-1).
\]
The selected values for \(2013\) are \(15,12,11,16\), so
\(t_{\max}=\min(14,11,10,15)=10\).  The eight-mask replay in
Table~\ref{tab:durer-masks} shows that no other admissible mask has selected
minimum \(11\), so this is also the unique maximin mask.  The full replay is
Claim IDs M24-C04, M24-C05, M24-C51, and M24-C64, certified in
\texttt{results/magic24\_certificate\_pack.json} and
\texttt{results/f2\_tesseract\_analysis.json}.
\end{proof}

\begin{table}[htbp]
\centering
\begin{tabular}{llrr}
\toprule
Mask & Selected values & \(t_{\max}\) & Terminal sum\\
\midrule
0231 & \(1,6,5,2\) & 0 & 34\\
0312 & \(1,9,11,3\) & 0 & 34\\
1203 & \(14,6,8,16\) & 5 & 29\\
1320 & \(14,9,10,13\) & 8 & 26\\
2013 & \(15,12,11,16\) & 10 & 24\\
2130 & \(15,7,5,13\) & 4 & 30\\
3021 & \(4,12,10,2\) & 1 & 33\\
3102 & \(4,7,8,3\) & 2 & 32\\
\bottomrule
\end{tabular}
\caption{The eight admissible one-incidence masks for the \Durer-complement
orientation used in this paper.}
\label{tab:durer-masks}
\end{table}

\begin{theorem}[Sagrada terminal as a finite retraction]
\label{thm:sagrada-retraction}
Index the cells by
\[
  x=(r_1,r_0,c_1,c_0)\in \F_2^4,
  \qquad r=2r_1+r_0,\quad c=2c_1+c_0.
\]
Let
\[
  M=\{0010,0100,1001,1111\}
\]
be the Sagrada mask plane and let \(u=1001\).  Define
\[
  F(x)=x+u\chi_M(x).
\]
Then the terminal square satisfies
\[
  D(10)_x=D_{F(x)}
\]
cell by cell.  Moreover \(F^2=F\), \(|\operatorname{im}F|=12\), and
\(M+u\) is the opposite one-incidence mask \(0231\).
\end{theorem}

\begin{proof}
Outside \(M\) the map \(F\) fixes the cell and \(D(10)=D\).  On the four
selected cells, the source values are \(15,12,11,16\), while the partner cells
in \(M+u\) have source values \(5,2,1,6\).  Subtracting \(10\) on \(M\) gives
exactly those partner values, hence \(D(10)_x=D_{F(x)}\).  The two planes
\(M\) and \(M+u\) are disjoint; therefore points of \(M\) move once, points of
\(M+u\) are fixed, and \(F\) is idempotent with four two-point fibers and eight
singleton fibers.  Thus \(|\operatorname{im}F|=12\).
\end{proof}

\begin{lemma}[Terminal duplicate-fiber direction]
\label{lem:terminal-duplicate-direction}
The direction \(u=1001\) is recoverable from the terminal square \(D(10)\)
itself: the duplicated terminal values are \(1,2,5,6\), and each duplicate
fiber has cell difference \(1001\).
\end{lemma}

\begin{proof}
The duplicate fibers are
\[
\begin{array}{c|c|c|c}
\text{value} & \text{first cell bits} & \text{second cell bits} & \text{xor}\\
\hline
1&0000&1001&1001\\
2&0100&1101&1001\\
5&0010&1011&1001\\
6&0110&1111&1001.
\end{array}
\]
All four duplicate fibers therefore determine the same direction.  This is a
terminal fact; it does not need the source square once \(D(10)\) is given.
\end{proof}

The closed bounded ray spectrum is
\[
\Ray^-_{D,M}=\{24,25,26,27,28,29,30,31,32,33,34\}.
\]
For the meet calculation, we use the strict downward part below the source
sum \(34\), namely \(\{24,\ldots,33\}\), because the meet compares Lo Shu
values above the base sum \(15\) with nontrivial downward perturbations below
the \Durer source.

\begin{lemma}[Terminal feasibility]
\label{lem:terminal-feasibility}
Let \(A\) be any normal order-four magic square, with entries
\(\{1,\ldots,16\}\), and let \(N\) be a one-incidence mask selecting four
cells.  If \(m\) is the smallest selected entry, then the terminal sum of the
bounded descent \(A-tN\) is \(35-m\).  Consequently, among the upper Lo Shu
sums
\[
\{18,21,24,27\},
\]
only \(24\) and \(27\) can occur as terminal sums of such a one-incidence
descent.
\end{lemma}

\begin{proof}
The source magic sum of a normal order-four square is \(34\).  A
one-incidence mask subtracts \(t\) once from each row, column, and main
diagonal, so the magic sum becomes \(34-t\).  The lower bound \(1\) allows
\[
t_{\max}=m-1,
\]
and therefore the terminal sum is \(34-(m-1)=35-m\).

To match an upper Lo Shu sum \(S\), one must have \(m=35-S\).  For
\(S=18\), this gives \(m=17\), impossible in \(\{1,\ldots,16\}\).  For
\(S=21\), this gives \(m=14\), but four distinct selected entries all at
least \(14\) cannot be chosen from \(\{1,\ldots,16\}\).  The remaining cases
\(S=24\) and \(S=27\) require \(m=11\) and \(m=8\), respectively, and are
not excluded by this terminal feasibility test.
\end{proof}

\section{The Strong Meet}
\label{sec:meet}

\begin{definition}[Weak and strong meet]
The weak meet is the numerical intersection of the upward Lo Shu spectrum
above \(15\) and the strict downward \Durer/Sagrada ray spectrum below \(34\).
The strong meet keeps only values that are non-degenerate on the Lo Shu side
and terminal on the \Durer/Sagrada side.
\end{definition}

\begin{theorem}[Strong bounded Lo Shu/\Durer-Sagrada meet]
\label{thm:strong-meet}
For the bounded categories above,
\[
\Meet_{\mathrm{weak}}=\{24,27\},
\qquad
\Meet_{\mathrm{strong}}=\{24\}.
\]
\end{theorem}

\begin{proof}
By Certified Proposition~\ref{prop:loshu-spectrum},
\[
\Spec^+_{3,9}(15)=\{18,21,24,27\}.
\]
By Certified Proposition~\ref{prop:sagrada-terminal}, the strict downward
\Durer/Sagrada ray contains
\[
\{24,25,26,27,28,29,30,31,32,33\}.
\]
The intersection is \(\{24,27\}\).  The value \(24\) is non-degenerate on the
Lo Shu side and terminal on the \Durer/Sagrada side.  The value \(27\) is the
all-9 Lo Shu fiber and occurs at the non-terminal \Durer ray time \(t=7\).
Therefore only \(24\) survives the strong meet rule.
\end{proof}

\begin{remark}
The value \(27\) is the control case.  It shows why the theorem needs the two
boundary conditions rather than a bare intersection of spectra.
\end{remark}

\begin{figure}[htbp]
\centering
\fbox{\begin{minipage}{0.88\linewidth}
\[
\begin{array}{ccl}
\text{Lo Shu upper spectrum}
  &:& 18\quad 21\quad \boxed{24}\quad 27_{\mathrm{deg}}\\[2mm]
\text{\Durer/Sagrada downward ray}
  &:& 33\quad 32\quad \cdots\quad 27\quad \cdots\quad \boxed{24}_{\mathrm{term}}
\end{array}
\]
\[
\Meet_{\mathrm{weak}}=\{24,27\},
\qquad
\Meet_{\mathrm{strong}}=\{24\}.
\]
\end{minipage}}
\caption{The strong meet keeps the value that is simultaneously
non-degenerate on the Lo Shu side and terminal on the \Durer/Sagrada side.}
\label{fig:strong-meet-spine}
\end{figure}

\section{Quaterne Transport}
\label{sec:quaternes}

For a \(4\times4\) square \(A\), let \(H_s(A)\) denote the family of
four-cell subsets whose entries sum to \(s\).  These subsets are not required
to be rows, columns, diagonals, or geometric lines.  They are the
four-uniform set-system layer of the square: the transport problem asks how
these labelled hyperedges move under a bounded mask perturbation.  We use
``quaterne'' for four-cell subsets, following the \Durer-pattern literature.
\Durer's \(34\)-sum quaternes are a classical theme in the literature on
melancholic magic squares \cite{trenklerMelancholic}.

\begin{certprop}[\Durer/Sagrada quaterne transport]
\label{prop:quaterne-transport}
The \Durer-complement square has
\[
|H_{34}(D)|=86.
\]
Against the Sagrada mask, those \(86\) quaternes have incidence distribution
\[
19,\;50,\;17
\]
for mask incidences \(0,1,2\).  The terminal square has
\[
|H_{24}(D(10))|=96,
\]
with source decomposition
\[
25+50+21,
\]
coming respectively from source sum \(24\) and incidence \(0\), source sum
\(34\) and incidence \(1\), and source sum \(44\) and incidence \(2\).
\end{certprop}

\begin{proof}[Certification]
This is Claim IDs M24-C07--M24-C09 and M24-C117, replayed in
\begin{itemize}[leftmargin=*]
\item \path{results/magic24_certificate_pack.json}
\item \path{results/inside_out_set_system_audit.json}
\end{itemize}
The transport rule is immediate from \(D(10)=D-10M\): a source quaterne with
Sagrada incidence \(k\) has terminal sum decreased by \(10k\).
\end{proof}

\begin{table}[htbp]
\centering
\begin{tabular}{lrrr}
\toprule
Source family & Mask incidence & Count & Terminal sum at \(t=10\)\\
\midrule
\(H_{34}(D)\) & 0 & 19 & 34\\
\(H_{34}(D)\) & 1 & 50 & 24\\
\(H_{34}(D)\) & 2 & 17 & 14\\
\midrule
Source sum 24 & 0 & 25 & 24\\
Source sum 34 & 1 & 50 & 24\\
Source sum 44 & 2 & 21 & 24\\
\bottomrule
\end{tabular}
\caption{The source \(H_{34}(D)\) incidence split and the terminal
\(H_{24}(D(10))\) decomposition.}
\label{tab:quaterne-transport}
\end{table}

\begin{certprop}[Terminal quotient count shadow]
\label{prop:terminal-count-shadow}
Let \(u=1001\) be the terminally recovered duplicate-fiber direction of
Lemma~\ref{lem:terminal-duplicate-direction}.  Among the \(96\) terminal
quaternes in \(H_{24}(D(10))\), exactly \(36\) are affine planes in
\(\F_2^4\) and \(60\) are non-affine.  If the non-affine quaternes are
quotiented by \(\langle u\rangle\), the resulting count vector is
\[
  60=32+12+16.
\]
More explicitly, the \(32\) mixed records have cell-xor \(u\), while the pure
non-affine records split modulo \(\langle u\rangle\) as \(4+8+16\), grouped as
\(12+16\) at the count-shadow level.
\end{certprop}

\begin{proof}[Certification]
This is replayed in
\path{results/sagrada_terminal_retraction_shadow.json} and summarized in
\path{results/SAGRADA_TERMINAL_RETRACTION_SHADOW_REPORT.md}.
The same replay also verifies Theorem~\ref{thm:sagrada-retraction} and
Lemma~\ref{lem:terminal-duplicate-direction}.  The statement is deliberately
only a count-shadow statement: it does not construct a direct record/profile
lift from \(D(10)\) to any endpoint-24 atlas class, and it does not identify
the non-normal terminal square \(D(10)\) with a normal order-four atlas record.
\end{proof}

\begin{remark}[Johnson-scheme reading]
The quaternes are \(4\)-subsets of the \(16\)-cell set, hence live in the
Johnson scheme \(J(16,4)\).  For a selected mask \(M\), write
\[
J_i(M)=\{Q:\ |Q|=4,\ |Q\cap M|=i\}.
\]
The transport identity is simply
\[
\sum_{c\in Q}(A-tM)_c=\sum_{c\in Q}A_c-t\,|Q\cap M|.
\]
The mask-only Phase-P replay in
\path{results/johnson_quaterne_layer_p.json} shows that this language
separates the main \(176=144+32\) selected-affine terminal-24 signature from
the remaining \(60\) outside-main records, but does not separate the \(144\)
exact-\(V_4\) records from the \(32\) selected-affine extras.  The anchored
replay in \path{results/johnson_anchor_refinement_p2.json} adds the fixed
translation \(V_4=\{0123,1032,2301,3210\}\) from the O5 layer as four
Johnson anchors; that anchored profile separates the \(144\), \(32\), and
\(60\) classes.  The direction-level replay in
\path{results/johnson_anchor_direction_p3.json} explains the exact
\(144\)-profile: relative to the fixed translation direction, the terminal
affine quaternes have the uniform inventory
\[
  1\ \text{same direction} \;+\; 5\ \text{line-intersection directions}
  \;+\; 3\ \text{complement directions},
\]
giving the anchored shadow
\[
  4(4,0,0,0)+20(2,2,0,0)+12(1,1,1,1).
\]
The follow-up replays in \path{results/johnson_followup_p4.json} and
\path{results/o5e_collision_forbidden_shadow_p5.json} sharpen this boundary.
For the exact \(144\), the \(36\) terminal-affine quaternes have inner
intersection distribution
\[
  |Q\cap Q'|=0:246,\qquad |Q\cap Q'|=1:192,\qquad |Q\cap Q'|=2:192,
\]
and every cell has degree \(9\).  More importantly, the \(32\)
selected-affine extras are exactly the records with positive forbidden
quotient-shadow mass.  If \(F_{2110}\) counts terminal-affine quaternes whose
anchored shadow is \((2,1,1,0)\), then
\[
  F_{2110}=0 \quad\text{for the exact }144,\qquad
  F_{2110}\in\{12,16\} \quad\text{for the }32\text{ extras}.
\]
The linear reason is short.  Let \(W_0\le F_2^4\) be the two-dimensional
direction of the fixed translation spread, and let \(A=a+L\) be a
domain-affine plane.  Under the quotient \(F_2^4\to F_2^4/W_0\), the image of
\(A\) is an affine subspace of dimension \(0\), \(1\), or \(2\); the nonempty
fibers of the restricted quotient map all have equal size.  Thus the sorted
spread-shadow of \(A\) is only
\[
  (4,0,0,0),\qquad (2,2,0,0),\qquad \text{or}\qquad (1,1,1,1),
\]
and \((2,1,1,0)\) is impossible for a domain-affine plane.  The P6 replay
\path{results/forbidden_shadow_split_p6.json} also shows that the
within-the-32 split is not independent: \(F_{2110}=12\) is exactly the
\(16\) two-diagonal translation-subset terminal records, while \(F_{2110}=16\)
is exactly the two \(8\)-record \(V_4\)-like terminal-set classes.  Thus the
\(144/32\) boundary is visible in O5-anchored Johnson data, not in the
mask-only Johnson strata.
\end{remark}

The \(50\) transported quaternes are the central bridge between the visible
\Durer quaterne layer and the terminal \(24\)-layer.  Section~\ref{sec:f2}
shows that \(36\) of these \(50\) transported quaternes are affine planes in
the \(F_2^4\) tesseract model.

\section{Permutation Diagonals and the Type-A Scope}
\label{sec:s4}

The \(24\) permutation diagonals of a \(4\times4\) array are indexed by
\(S_4\).  This is also the chamber set of the type-A Coxeter complex \(A_3\).
In the present paper the type-A language is used only at this finite
chamber-indexing level.

\begin{certprop}[Permutation-diagonal subgroup drop]
\label{prop:d4-v4}
In the source \Durer-complement square, the target permutation diagonals of
sum \(34\) are
\[
\{0123,0213,1032,1302,2031,2301,3120,3210\},
\]
an order-\(8\) subgroup isomorphic to \(D_4\).  For every \(t>0\) along the
Sagrada ray, the target diagonals of the current magic sum \(34-t\) are
\[
\{0123,1032,2301,3210\},
\]
the Klein four subgroup \(V_4\).
\end{certprop}

\begin{proof}[Certification]
This is Claim IDs M24-C10 and M24-C11 in the crosswalk, replayed in
\texttt{results/magic24\_certificate\_pack.json}.
\end{proof}

\begin{table}[htbp]
\centering
\begin{tabular}{lll}
\toprule
Ray position & Target diagonal set & Group type\\
\midrule
\(t=0\) & \(0123,0213,1032,1302,2031,2301,3120,3210\) & \(D_4\)\\
\(t>0\) & \(0123,1032,2301,3210\) & \(V_4\)\\
\bottomrule
\end{tabular}
\caption{The \Durer/Sagrada permutation-diagonal drop.}
\label{tab:d4-v4}
\end{table}

\begin{certprop}[Type-A guardrail]
\label{prop:type-a-guardrail}
In the \(A_3/S_4\) chamber model, these \(D_4\) and \(V_4\) subsets are
subgroup/coset chamber tilers, not standard poset cones \(L(P)\).  Their
common-halfspace closure is all of \(S_4\), and their induced chamber graphs
are disconnected.
\end{certprop}

\begin{proof}[Certification]
This is Claim IDs M24-C12, M24-C41--M24-C43, and M24-C49, replayed in
\begin{itemize}[leftmargin=*]
\item \path{docs/TYPE_A_SUBGROUP_TILERS.md}
\item \path{docs/TYPE_A_POSET_CONE_COMPARISON.md}
\item \path{results/s4_chamber_fingerprints.json}
\item \path{results/s4_poset_cone_comparison.json}
\end{itemize}
\end{proof}

\begin{remark}
The general subgroup/coset chamber-tiler component theorem belongs to the
separate Type-A/Tetra paper line.  In this paper \(D_4\to V_4\) is used only
as a \Durer-side enrichment and as a finite guardrail: subgroup tilers are not
automatically poset cones.
\end{remark}

\section{The Binary Tesseract Layer}
\label{sec:f2}

Sudbery's tesseract reading of \Durer's square suggests labelling the \(16\)
cells by \(D_{ij}-1\), hence by the elements of \(\F_2^4\)
\cite{sudberyTesseract}.  In the orientation used here this is not merely a
set-theoretic labelling.  It is linear in row and column bits.

\begin{certprop}[Linear tesseract coordinatization]
\label{prop:f2-linear}
Write the row index and column index as two-bit vectors
\[
r=(r_0,r_1),\qquad c=(c_0,c_1).
\]
If
\[
D_{r,c}-1=(\ell_0,\ell_1,\ell_2,\ell_3)\in \F_2^4,
\]
then
\[
\ell_0=r_1+c_0+c_1,\quad
\ell_1=r_0+c_0+c_1,\quad
\ell_2=r_0+r_1+c_0,\quad
\ell_3=r_0+r_1+c_1,
\]
with all operations in \(\F_2\).  Thus the cell-value map is a rank-\(4\)
linear automorphism of \(\F_2^2\oplus \F_2^2\).  Consequently all \(24\)
permutation diagonals are affine planes in the \Durer \(F_2^4\) model.
\end{certprop}

\begin{proof}[Certification]
This is Claim IDs M24-C73 and M24-C74, replayed in
\begin{itemize}[leftmargin=*]
\item \path{docs/F2_TESSERACT_LAYER.md}
\item \path{results/f2_tesseract_analysis.json}
\item \path{tests/test_f2_tesseract.py}
\end{itemize}
The final assertion also uses
\(\AGL(2,2)\cong S_4\): every permutation of the four column labels is
affine over \(\F_2^2\), so each graph \(x\mapsto \sigma(x)\) is an affine
plane.
\end{proof}

\begin{certprop}[Affine source and terminal transport]
\label{prop:f2-transport}
There are \(140=35\cdot4\) affine planes in \(\F_2^4\).  Inside
\(H_{34}(D)\), the affine quaternes are exactly
\[
13\text{ balanced directions}\times4\text{ cosets}=52.
\]
The Sagrada mask has direction \(W\).  Among the \(13\) balanced directions,
\(9\) are complementary to \(W\), and these contribute
\[
9\cdot4=36
\]
terminal affine quaternes at \(D(10)\).  These \(36\) are precisely the pure
transport layer from source sum \(34\) and Sagrada incidence \(1\).
\end{certprop}

\begin{proof}[Certification]
This is Claim IDs M24-C65--M24-C67 and M24-C75--M24-C77.  The replay is in
\texttt{docs/F2\_TESSERACT\_LAYER.md} and
\texttt{results/f2\_tesseract\_analysis.json}.  A direction is balanced
exactly when no coordinate bit is constant on it; these are the \(13\)
directions whose cosets have integer value sum \(34\).  Complementarity to
the Sagrada direction gives one Sagrada point in each coset, hence terminal
sum \(24\) after subtracting \(10\).
\end{proof}

\begin{table}[htbp]
\centering
\begin{tabular}{lr}
\toprule
\(F_2^4\) layer & Count\\
\midrule
Ambient affine planes & \(35\cdot4=140\)\\
Balanced source directions & 13\\
Affine source quaternes in \(H_{34}(D)\) & \(13\cdot4=52\)\\
Balanced directions complementary to Sagrada & 9\\
Terminal affine quaternes in \(H_{24}(D(10))\) & \(9\cdot4=36\)\\
\bottomrule
\end{tabular}
\caption{Finite-affine summary of the tesseract transport layer.}
\label{tab:f2-transport}
\end{table}

The same \(F_2^2\) language gives a concise interpretation of
Proposition~\ref{prop:d4-v4}:
\[
D_4=\{x\mapsto x+b\}\cup\{x\mapsto \operatorname{swap}(x)+b\},
\qquad
V_4=\{x\mapsto x+b\}.
\]
The ray collapses the target diagonal family from translations plus
bit-swap translations to translations alone.

\begin{remark}
The \(F_2^4\) layer enriches the \Durer/Sagrada side of the meet.  It does
not make \(24\) universal, and in the full order-four atlas it tracks the
exact canonical \(V_4\) affine subclass rather than only the historical
\Durer/Sagrada cell.
\end{remark}

\section{Parity-Plane Mechanism}
\label{sec:parity-plane}

The \(F_2^4\) layer also gives a local explanation of why the endpoint-24
condition is visible at the \Durer side of the meet.  This explanation is
scoped: it does not make \(24\) universal, and it does not replace the
order-four endpoint atlas.  It identifies one parity-plane obstruction inside
the normalized binary model used in this paper.

Throughout this section, \(L(x)\) denotes the zero-based four-bit label at
cell \(x\).  The local residual reduction supplies a visible parity coefficient
\(B_{11}\), defined operationally in the public replay by the identity
\[
\bigoplus_{x\in\Pi_{NW}}L(x)=B_{11}.
\]
The six residual values below are the nonzero cases left by that normalized
reduction; this section checks those cases rather than deriving a table-free
classification of all residual forms.

\begin{definition}[NW parity plane and bit-balance]
Use zero-based row and column coordinates and write a cell as
\(x=(r_1,r_0,c_1,c_0)\in\F_2^4\).  The NW parity plane is the even-row,
even-column four-cell set, equivalently the row-major index set
\[
\Pi_{NW}=\{0,2,8,10\}.
\]
For zero-based labels in \(\{0,\ldots,15\}\), a four-label set is
bit-balanced if each of the four coordinate bits occurs exactly twice.  Such a
set has integer label sum
\[
2(1+2+4+8)=30.
\]
\end{definition}

\begin{lemma}[Affine-plane sum]
\label{lem:affine-plane-sum}
Let \(u,w\in\F_2^4\) be nonzero distinct directions, and let
\[
P=\{o,o\oplus u,o\oplus w,o\oplus u\oplus w\}\subset \F_2^4.
\]
Then \(P\) has four cells, and the zero-based integer label sum of \(P\) is \(30\) if and only if the two
directions cover all four bits:
\[
u\mathbin{\mathrm{OR}} w=15.
\]
\end{lemma}

\begin{proof}
In a fixed bit position, the four labels on \(P\) are balanced exactly when at
least one of \(u\) and \(w\) moves that bit.  If both direction bits are zero,
that bit is constant on \(P\).  Thus all four bits occur twice exactly when
\(u\) and \(w\) cover all four bit positions, which is the displayed OR
condition.
\end{proof}

\begin{certprop}[Parity-plane obstruction, finite residual form]
\label{prop:parity-plane-mechanism}
In the normalized residual split used here, the public certificate checks the
nonzero residual forms supplied by the parity-plane reduction and excludes each
of them from carrying the NW parity-plane bit-balance.  The visible branch
\(B_{11}\ne0\) is excluded by
\[
\bigoplus_{x\in\Pi_{NW}} L(x)=B_{11},
\]
and the remaining nonzero residual forms checked in the replay are the two
transpose-symmetric row/column patterns
\[
(0,0,v,v),\qquad (0,v,0,v),
\qquad v\in\{2,4,6,8,12,14\}.
\]
For these six values the residual direction constraints force the NW
directions to miss a coordinate bit, so the NW parity plane is not
bit-balanced.
\end{certprop}

\begin{proof}[Certification and finite residual check]
The identity \(\bigoplus_{x\in\Pi_{NW}}L(x)=B_{11}\) gives the first split:
bit-balance implies xor zero, so \(B_{11}\ne0\) is impossible.  The public
finite check then handles the nonzero residual branch with \(B_{11}=0\).  For
the row residual, write the affine anchor as
\[
A(x)=o\oplus a r_1\oplus b r_0\oplus g c_1\oplus d c_0.
\]
The row and column plane equations, the two diagonal equations, and the
bijective-toggle condition leave the six values shown above.  In every
surviving row case the replay verifies \(a=v\), and every admissible \(g\) has
zero in the critical bit \(\nu_2(v/2)\).  Hence the NW directions \(a,g\) do
not cover all four bits:
\[
a\mathbin{\mathrm{OR}}g\ne15.
\]
By Lemma~\ref{lem:affine-plane-sum}, the NW plane is not bit-balanced.  The
column pattern follows by transposing row and column variables.

The certification artifacts are
\begin{itemize}[leftmargin=*]
\item \path{docs/PARITY_PLANE_MECHANISM.md},
\item \path{results/parity_plane_mechanism.json},
\item \path{results/PARITY_PLANE_MECHANISM_REPORT.md},
\item \path{tests/test_parity_plane_mechanism.py}.
\end{itemize}
The last step is intentionally recorded as a six-value finite residual check;
we do not claim here a table-free symbolic proof of the residual obstruction.
\end{proof}

\begin{remark}[Scope]
Proposition~\ref{prop:parity-plane-mechanism} is a mechanism for the
Shu--\Durer bridge inside this normalized binary model.  It is not a universal
statement about endpoint \(24\) across all magic-square settings.
\end{remark}

\section{Bounded-Polytope Guardrail}
\label{sec:polytope}

Magic squares can also be studied as lattice points in polyhedral cones and
bounded polytopes
\cite{wardMagicVectorSpaces,ahmedMagicSquares,ahmedDeLoeraHemmecke}.  Here
that language serves mainly as a guardrail.

\begin{certprop}[The terminal square is not a vertex]
\label{prop:not-vertex}
The terminal square \(D(10)\) lies on the boundary of the bounded order-four
magic-square polytope, but it is not a vertex of either the free-sum bounded
polytope or the fixed-sum \(S=24\) bounded polytope.  The full fixed-sum
\(S=24\) bounded polytope has affine dimension \(7\) and \(292\) vertices,
with \(180\) integral and \(112\) semi-integral vertices.  The point \(D(10)\)
has a two-vertex barycentric certificate
\[
D(10)=\frac45 V_{126}+\frac15 V_{140}
\]
in the full vertex catalogue.
\end{certprop}

\begin{proof}[Certification]
This is Claim IDs M24-C52, M24-C56, M24-C60, M24-C61, and M24-C63, replayed
in the bounded-polytope and fixed-sum audit artifacts:
\begin{itemize}[leftmargin=*]
\item \path{docs/BOUNDED_MAGIC_POLYTOPE.md}
\item \path{docs/FIXED_SUM24_POLYTOPE_AUDIT.md}
\item \path{results/bounded_magic_polytope.json}
\item \path{results/fixed_sum24_polytope_audit.json}
\end{itemize}
\end{proof}

Thus the \(4\times4\) side is a bounded ray endpoint, not a shared vertex
mechanism with the Lo Shu side.  This negative statement prevents the paper
from replacing one overclaim by another.

\section{Order-Four Extension}
\label{sec:order4}

The \Durer/Sagrada endpoint is not unique in the wider normal order-four
dataset.  The companion replay first reconstructs the classical counts
\[
7040
\]
raw normal order-four magic squares and
\[
880
\]
essential representatives \cite{ollerenshawBondi}.

\begin{certprop}[Endpoint-24 atlas and exact \(V_4\) subclass]
\label{prop:order4-extension}
Across the \(880\) essential order-four representatives and the eight
admissible one-incidence masks, endpoint \(24\) occurs in \(236\)
essential square-mask pairs.  In those terminal-24 records, affine cell-value
labelling occurs in exactly \(144\) pairs, and this equals the exact
canonical \(V_4\) terminal subclass in the chosen essential orientation.  The
main inside-out terminal signature is a broader \(176\)-pair class:
\[
176=144+32.
\]
\end{certprop}

\begin{proof}[Certification]
This is Claim IDs M24-C15--M24-C17, M24-C78--M24-C80, M24-C124--M24-C127,
and M24-C139.  The replay artifacts are
\begin{itemize}[leftmargin=*]
\item \path{data/order4_normal_essential_880.json}
\item \path{results/order4_endpoint_spectrum.json}
\item \path{results/order4_f2_extension.json}
\item \path{results/exact_v4_affine_class_audit.json}
\item \path{results/inside_out_main_signature_split.json}
\end{itemize}
\end{proof}

The conclusion is qualitative rather than uniqueness-based.  \Durer/Sagrada is
a historical canonical representative inside a larger structured terminal-24
landscape.  The \(32\) extras in the \(176=144+32\) signature are structured
near misses: earlier finite invariants did not explain them, while the
current O5-anchored Johnson refinement separates them from the exact
\(V_4\) class by positive forbidden quotient-shadow mass.  Their broader
conceptual role remains a follow-up question rather than a strong-meet
mechanism.

\begin{figure}[htbp]
\centering
\fbox{\begin{minipage}{0.86\linewidth}
\[
\begin{array}{rcl}
236 & = & 176 \sqcup 60
      \quad\text{terminal-24 records split into main and outside-main}\\
176 & = & 144 \sqcup 32
      \quad\text{main class split into exact }V_4\text{ and extras}\\
144 & \supset & 8
      \quad\text{the historical \Durer/Sagrada quotient cell.}
\end{array}
\]
\end{minipage}}
\caption{Order-four terminal-24 containment diagram.  The most stable finite
class detected by the atlas is the \(144\)-pair globally affine exact-\(V_4\)
subclass; the historical
\Durer/Sagrada cell is a small canonical representative inside it.}
\label{fig:order4-containment}
\end{figure}

\section{Certificate Architecture}
\label{sec:certificates}

The main body uses a narrow claim set.  Table~\ref{tab:claim-map} summarizes
the claim levels and keeps the computational layers in their proper roles.

\begin{table}[htbp]
\centering
\begin{tabular}{p{0.25\linewidth}p{0.25\linewidth}p{0.38\linewidth}}
\toprule
Layer & Paper role & Primary repository control\\
\midrule
Strong meet & theorem & \texttt{CLAIM\_CROSSWALK.md}\\
Lo Shu and Sagrada ray & certified propositions & certificate pack\\
Quaterne transport & certified proposition & certificate pack, inside-out audit\\
Sagrada terminal retraction/count shadow & theorem/certprop & terminal retraction replay\\
\(D_4\to V_4\), Type-A scope & certified propositions & \(S_4\) reports\\
\(F_2^4\), bounded polytope & certified propositions / guardrails & Phase F/E reports\\
Parity plane & certified finite check & parity-plane report/test\\
880/236/144/176 extension & certified proposition / appendix & Phase C/H/O reports\\
APD, Markov, Hilbert, \(N\)-poset & appendix audits & reproducibility package\\
\bottomrule
\end{tabular}
\caption{Paper-v2 claim levels and repository controls.}
\label{tab:claim-map}
\end{table}

The strongest methodological lesson is that finite certificates are most
useful when they also state what they do \emph{not} prove.  This is why the
APD/PTE, Morse-Hedlund, Markov, Hilbert, and pointed five-state analyses
remain appendices or follow-up problems.

\appendix

\section{\texorpdfstring{Endpoint Spectrum for the \(880\)-Representative Atlas}{Endpoint Spectrum for the 880-Representative Atlas}}
\label{app:endpoint-atlas}

The endpoint spectrum across the \(880\) essential representatives and the
eight one-incidence masks is
\[
\begin{array}{c|rrrrrrrrrrrrr}
\text{endpoint} & 22&23&24&25&26&27&28&29&30&31&32&33&34\\
\hline
\text{count} & 176&44&236&268&676&108&328&376&772&420&884&992&1760
\end{array}
\]
Endpoint \(24\) is therefore not rare in the strongest possible sense.  Its
role in this paper comes from the strong meet and from the \Durer-side
structure, not from absolute uniqueness in the order-four atlas.

The enriched filter ladder
\[
7040\to236\to220\to160\to152\to144
\]
passes through endpoint \(24\), subgroup terminal diagonals, Klein-four
profile, terminal quaterne count \(96\), selected values
\(\{11,12,15,16\}\), and exact canonical \(V_4\).

\section{\texorpdfstring{The Exact Canonical \(V_4\) Subclass and the \Durer/Sagrada Cell}{The Exact Canonical V4 Subclass and the Durer/Sagrada Cell}}
\label{app:v4-subclass}

The \(144\)-pair exact canonical \(V_4\) subclass has selected values
\[
(11,12,15,16),
\]
terminal diagonal set
\[
\{0123,1032,2301,3210\},
\]
and terminal quaterne decomposition \(25+50+21\).

The Phase-O structural audit sharpens this class inside the \(236\)
terminal-24 records.  The following conditions are equivalent there:
\[
\begin{gathered}
\text{exact canonical }V_4,\\
\text{global affine cell-value labelling},\\
\text{zero affine defect},\\
\text{full translation terminal }V_4,\\
\text{preservation of all }140\text{ domain affine planes},\\
\text{preservation of all }24\text{ permutation-diagonal affine planes}.
\end{gathered}
\]
The selected-mask affine condition alone remains too broad: it gives the
\(176=144+32\) main class.  The \(32\) selected-affine extras are the first
near-miss layer, with affine defect \(4\), \(76/140\) domain affine planes
preserved, and \(8/24\) permutation-diagonal affine planes preserved.

There is also a wider affine-layer explanation of the number \(144\).  It is
useful to state it separately, because it no longer depends on filtering the
terminal-24 atlas.

\begin{lemma}[Balanced directions]
\label{lem:balanced-directions}
In \(\F_2^4\), exactly \(13\) two-dimensional linear subspaces are balanced,
meaning that no coordinate bit is constant on the subspace.
\end{lemma}

\begin{proof}
There are
\[
\left[\begin{matrix}4\\2\end{matrix}\right]_2=35
\]
two-dimensional subspaces of \(\F_2^4\).  A direction is bad exactly when it
is contained in at least one coordinate hyperplane.  Each coordinate
hyperplane is \(\F_2^3\) and contains
\[
\left[\begin{matrix}3\\2\end{matrix}\right]_2=7
\]
two-dimensional subspaces.  Two coordinate hyperplanes intersect in one
two-dimensional coordinate plane, while three coordinate hyperplanes
intersect in a line and contain no two-dimensional subspace.  Hence
\[
\#\{\text{bad directions}\}=4\cdot7-\binom42=22,
\]
so \(35-22=13\) directions are balanced.
\end{proof}

\begin{certprop}[Affine normal count]
\label{prop:affine-normal-count}
There are \(432\) essential normal order-four magic squares whose cell-value
labelling \(x\mapsto Q_x-1\) is globally affine over \(\F_2^4\).  Moreover,
across the eight admissible one-incidence masks, these representatives give
\[
432\cdot8=3456
\]
affine square-mask pairs, uniformly distributed as
\[
24\text{ selected value planes}\times144\text{ pairs}.
\]
\end{certprop}

\begin{proof}[Certification and finite-linear derivation]
Write a globally affine normal square as
\[
Q_x-1=A x+b,\qquad A\in GL(4,2),\quad b\in\F_2^4.
\]
If a four-cell affine plane has balanced direction, then every coordinate bit
is \(1\) on exactly two of its four labels.  Therefore its label sum is
\[
2(8+4+2+1)=30,
\]
and its value sum is \(30+4=34\).  Thus a linear part \(A\) is good exactly
when it sends the row, column, and diagonal directions of the cell domain to
balanced directions in value space; the main and anti-diagonal are parallel
affine planes and share the same direction.

The finite \(F_2\)-linear replay gives the following count.  The
complementarity graph on the \(13\) balanced directions has \(6\) triangles,
so there are \(6\cdot6=36\) ordered pairwise-complementary balanced image
triples.  For each ordered image triple there are \(6\) linear maps, hence
\[
36\cdot6=216
\]
good linear parts.  Offsets are free, giving
\[
216\cdot16=3456
\]
raw affine normal squares.  Since a normal square has all entries distinct,
no nontrivial square symmetry fixes it cellwise; hence each square-symmetry
orbit has size \(8\), and
\[
3456/8=432
\]
essential affine representatives.

The same replay shows that each affine representative sends the eight
admissible masks to exactly
\[
8=2\text{ coordinate directions}\times4\text{ cosets}
\]
selected value planes.  The two directions form one of the three matchings of
the four coordinate axes:
\[
1,2\,|\,4,8,\qquad 1,4\,|\,2,8,\qquad 1,8\,|\,2,4.
\]
The \(24\) value-bit permutations preserve the affine normal layer and induce
the natural \(S_4\to S_3\) action on these three matchings.  Each matching has
stabilizer size \(8\), so the \(432\) affine representatives split as
\[
144+144+144.
\]
This yields the uniform \(24\times144\) selected-value-plane split.

The replay artifacts are
\begin{itemize}[leftmargin=*]
\item \path{results/affine_normal_layer.json}
\item \path{results/affine_normal_spread_structure.json}
\item \path{results/affine_normal_count_derivation.json}
\end{itemize}
\end{proof}

\begin{lemma}[Admissible-mask torsor]
\label{lem:mask-torsor}
The square-symmetry group of order \(8\) acts freely and transitively on the
eight admissible one-incidence masks.
\end{lemma}

\begin{proof}[Certification]
The eight masks are
\[
0231,\;0312,\;1203,\;1320,\;2013,\;2130,\;3021,\;3102.
\]
The square-symmetry replay has one orbit of size \(8\) and trivial stabilizer
for each mask.  This is recorded in
\(\texttt{results/affine\_normal\_count\_derivation.json}\).
\end{proof}

The two appearances of \(3456\) count different but compatible sets:
\[
\begin{array}{rcl}
216\cdot16&=&3456\quad\text{raw affine normal squares},\\
432\cdot8&=&3456\quad\text{essential affine square-mask pairs}.
\end{array}
\]
Lemma~\ref{lem:mask-torsor} explains why these counts mirror each other after
passing between oriented squares and essential representatives with a chosen
admissible mask.

Finally, if a selected value plane is \(P\), then its bounded downward endpoint
is
\[
34-(\min P-1)=35-\min P.
\]
Among the \(24\) selected value planes in Proposition~\ref{prop:affine-normal-count},
the unique plane with endpoint \(24\) is
\[
\{11,12,15,16\}.
\]
Every record in this plane has full translation terminal \(V_4\).  Thus the
`\(144\)' is not only a terminal-24 census count: it is one uniform value
plane inside the globally affine normal order-four layer.

\begin{certprop}[Affine translation torsors]
\label{prop:affine-translation-torsors}
For every one of the \(3456\) globally affine square-mask pairs, the four
fixed translation diagonals
\[
0123,\qquad 1032,\qquad 2301,\qquad 3210
\]
map to the four cosets of one balanced value-space direction \(W\).  If
\(P=a+U\) is the selected value plane, then
\[
\F_2^4=U\oplus W,
\]
so \(P\) intersects each \(W\)-coset in exactly one point.  The resulting map
from the four translation labels to the four points of \(P\) is an affine
bijection.

Moreover, this structure is uniform over all selected value planes.  For
each selected direction \(U\), let \(Q\) be its coordinate partner in the
matching of coordinate axes.  The nine balanced complements to \(U\) are
graphs of linear maps
\[
Q\longrightarrow U
\]
with no zero rows.  They split as
\[
9=6+3,
\]
where the six used directions are precisely the invertible maps
\(\GL(2,2)\), and the three unused directions are rank-one graphs.  Hence
each selected plane has
\[
6\text{ used directions}\times24\text{ records}=144
\]
affine square-mask pairs.
\end{certprop}

\begin{proof}[Certification]
This is a finite replay of the Phase-O5 torsor refinement.  The kernel
guardrail is recorded in
\(\texttt{results/kernel\_v4\_o5.json}\): the value-bit kernel \(V_4\) of
the \(S_4\to S_3\) action on coordinate matchings is not the terminal
translation \(V_4\).  The quotient-section statement is replayed in
\path{results/terminal_parallel_torsor_o5b.json}, where all \(3456\)
affine square-mask pairs have balanced translation-image
direction, complementary selected direction, intersection pattern
\((1,1,1,1)\), and affine translation-to-selected point map.

The six-of-nine split for the terminal-24 plane is replayed in
\(\texttt{results/six\_of\_nine\_o5c.json}\), and the all-selected-plane
generalization is replayed in
\path{results/selected_plane_six_of_nine_o5d.json}.  There are \(24\)
selected planes; each has \(144\) records, nine balanced
complements to its selected direction, six invertible graph complements,
three rank-one graph complements, and zero violations of the rule that the
used directions are exactly the six invertible graph complements.
\end{proof}

\begin{remark}[The final \(24\)-factor]
The count
\[
3456=24\text{ selected planes}\times6\text{ invertible directions}
\times24\text{ records}
\]
does not mean that every \((P,W)\)-fiber is the full
\(\AGL(2,2)\)-torsor of affine point maps.  The Phase-O5e/O5f audit shows
that only \(96\) of the \(144\) \((P,W)\)-fibers realize all \(24\) affine
point maps; \(48\) fibers have duplicated point maps.  The safe
parametrization is mask-refined: \((P,W,\phi,\mathrm{mask})\) is injective,
with \(24\) distinct \((\phi,\mathrm{mask})\) pairs in every \((P,W)\)-fiber.
The P5 replay sharpens the negative statement: in every \((P,W)\)-fiber the
point-map collision graph is a matching with \(0\), \(4\), or \(8\) collision
edges, and the pair consisting of selected-plane direction and translation
direction determines which case occurs.  The replays are
\path{results/torsor_parametrization_o5e_o5f.json} and
\path{results/o5e_collision_forbidden_shadow_p5.json}.
\end{remark}

\begin{table}[htbp]
\centering
\begin{tabular}{lrl}
\toprule
Source type & Count & Source diagonal set\\
\midrule
S1 & 96 & \(\{0123,1032,2301,3210\}\)\\
S2 & 16 & \(\{0123,0132,1023,1032,2301,2310,3201,3210\}\)\\
S3 & 16 & \(\{0123,0213,1032,1302,2031,2301,3120,3210\}\)\\
S4 & 16 & \(\{0123,0321,1032,1230,2103,2301,3012,3210\}\)\\
\bottomrule
\end{tabular}
\caption{Source diagonal types inside the \(144\)-pair exact canonical
\(V_4\) subclass.  The \Durer/Sagrada representative lies in type S3.}
\label{tab:v4-source-types}
\end{table}

The \Durer/Sagrada quotient cell has \(8\) pairs.  In the current finite atlas
it is characterized by APD vector \((0,0,0,55296)\), source type S3, exact
terminal \(V_4\), and terminal quaterne decomposition \(25+50+21\).  The
essential square indices are
\[
112,\;174,\;289,\;359,\;789,\;802,\;834,\;849.
\]
All \(8\) records are associated and complement-fixed after
canonicalization.  The canonical \Durer-complement representative is square
index \(174\); under canonicalization the Sagrada mask appears as \(1203\),
with selected values \(12,11,15,16\).

\section{Inside-Out, APD, and Permutation-Polytope Audits}
\label{app:secondary-audits}

The following audits support the paper while staying outside the main theorem
mechanism.

\begin{itemize}[leftmargin=*]
\item \textbf{Inside-out set systems.}  The row, column, diagonal, mask,
source-quaterne, and terminal-quaterne incidence systems have exact
rank/\(\SNF\)/kernel and colored-automorphism replays.  The main terminal
signature is the broad \(176=144+32\) class; the later O5-anchored Johnson
replay separates its \(144\) exact-\(V_4\) records from the \(32\) extras by
the forbidden quotient-shadow count \(F_{2110}\).
\item \textbf{APD/PTE.}  Along the Sagrada ray,
\[
\APD_1(t)=0,\qquad \APD_2(t)=0,\qquad
\APD_3(t)=-24t(t-4)(t-16).
\]
Moreover \(\APD_3(D(t))=6\det \widehat E(t)\) in the centered local model.
This gives a useful parity-power or Prouhet--Tarry--Escott reading of the
ray \cite{takemuraAPD,takemura880}, but endpoint \(24\) is selected by
bounded terminality, not by APD vanishing.
\item \textbf{Morse-Hedlund/Prouhet family fit.}  The \Durer-complement source
fits the displayed Morse-Hedlund/Prouhet order-four family, but the Sagrada
direction, terminal point, and full ray do not under the tested transform
scope \cite{sergeyevSymmetricSums}.
\item \textbf{Permutation polytopes.}  The subgroup polytopes
\[
P(D_4)=\conv\{P_\sigma:\sigma\in D_4\},
\qquad
P(V_4)=\conv\{P_\sigma:\sigma\in V_4\}
\]
fit naturally into the permutation-polytope language
\cite{baumeisterPermutationPolytopes,guralnickPerkinson}.  The face audit
shows that \(P(V_4)\) is a tetrahedron, but it is not a face of \(P(D_4)\) or
of the Birkhoff polytope \(B_4\).
\end{itemize}

\section{Markov and Hilbert Follow-Up Audits}
\label{app:markov-hilbert}

The terminal-24 records also support two technical follow-up directions.

\begin{itemize}[leftmargin=*]
\item The first Markov-style graph uses primitive \(\{-1,0,1\}\) fixed-sum
diagonal-magic kernel moves.  It has \(236\) nodes, \(247\) edges, \(98\)
connected components, and \(50\) isolated vertices.  This is a small-move
graph, not a complete Markov-basis theorem.
\item The finite Hilbert-style audit through magic sum \(8\) finds \(20\)
checked atoms:
\[
8\text{ at sum }1,\qquad 12\text{ at sum }2.
\]
All \(236\) terminal-24 endpoints decompose in this checked atom set.  This
is not yet a proof of the full Hilbert basis of the nonnegative
diagonal-magic cone.
\end{itemize}

\section{The Pointed Five-State Bridge}
\label{app:five-state}

The Lo Shu \(S=24\) fiber has five points.  A separate audit compares this
five-state object with the graph of linear extensions of an \(N\)-poset and
with the order-ideal lattice \(J(A_2^+)\), using standard poset language
\cite{naatzLinearExtensions,stanleyPosetPolytopes}.

The natural Lo Shu adjacency is a diamond with a central point and four
boundary points, so it is not directly the \(N\)-poset linear-extension
graph.  After adding a boundary cycle on the four nonconstant points and
marking one boundary point, however, the resulting pointed graph is a
4-cycle with one pendant vertex.  It is isomorphic to the \(N\)-poset
linear-extension graph and to the Hasse graph of \(J(A_2^+)\).

The bridge is controlled but noncanonical: every boundary mark works, and no
current \Durer-side structure selects one.  It remains an open problem rather
than a foundation for the meet theorem.

\section{Open Problems}
\label{app:open}

The most useful follow-up questions are:
\begin{enumerate}[leftmargin=*]
\item Extend the finite-linear Phase-O/O2/O3/O4/O5 explanation of the exact
\(V_4\) affine class beyond the normal order-four affine layer, or show that
this layer is the natural boundary of the phenomenon.
\item Generalize the forbidden quotient-shadow obstruction for the \(32\)
structured extras inside the \(176=144+32\) inside-out signature, or show that
it is specific to this fixed order-four selected-affine terminal layer.
\item Develop subgroup/coset chamber tilers as a sibling theory to type-A
poset-cone tilers.  This belongs primarily to the Type-A/Tetra paper line;
Magic supplies the \(D_4\to V_4\) seed example.
\item Decide whether a complete Markov-basis or Hilbert-basis treatment is
worth a separate technical paper.
\item Find, or rule out, a natural \Durer-to-Lo-Shu boundary selector for the
pointed five-state bridge.
\end{enumerate}

\bibliographystyle{alpha}
\bibliography{refs}

\end{document}
